% !TeX root = ../../Book3.tex
% 纲要标题：### 13.4 完备化的局部性质
% 纲要位置：工作记录/计划纲要.md，第 2229 行。
% 纲要细目（写作范围）：
% * 完备化与张量积
% * Noether 局部环完备化的忠实平坦性
% * 极大理想、剩余域及维数
% * 证明顺序固定为：有限模完备化正合性、\(M\otimes_R\widehat R\cong\widehat M\)、第12.1节的理想张量与局部判据、平坦性、剩余域非零及第12.5节的忠实平坦性判据
% * 极大理想进完备化的维数比较使用第11章的伴随分次与 Hilbert–Samuel 理论；有限展示模的下降仍是独立选读，不作为此处证明前提


\section{完备化的局部性质}
\label{b3:sec:1304}


上一节已经证明完备化保持有限模短正合列，并把有限模的完备化写成张量积。本节才开始讨论平坦性；这样，平坦性是前面两个结果的推论，而不是完备化正合性的假设。

\subsection{完备化与张量积}
\label{b3:subsec:1304:01}

对有限模 \(M\)，\cref{b3:thm:completion-tensor}给出的同构与模同态相容：若 \(f:M\rightarrow N\)，两边都把 \(a\otimes m\) 送到 \(a\,\iota_N\Bigl(f(m)\Bigr)\)。因此可以把有限模之间的完备化映射统一视为沿 \(R\rightarrow\widehat R\) 的标量扩张。

\begin{proposition}\label{b3:prop:completion-finite-tensor-products}
设 \(R\) 为交换 Noether 环、\(I\subseteq R\) 为理想，\(M,N\) 为有限生成 \(R\) 模；记 \(\widehat R\)、\(\widehat M\)、\(\widehat N\) 及 \(\widehat{M\otimes_RN}\) 为各自的 \(I\) 进完备化。则
\[
\widehat{M\otimes_RN}
\cong\widehat M\otimes_{\widehat R}\widehat N.
\]
\end{proposition}
\begin{proof}
左边由张量表示等于 \(\widehat R\otimes_R(M\otimes_RN)\)。右边把两因子写成 \(\widehat R\otimes_RM\)、\(\widehat R\otimes_RN\)，映射
\[
(a\otimes m)\otimes(b\otimes n)\longmapsto ab\otimes(m\otimes n)
\]
保持各个平衡关系，逆映射为
\(c\otimes(m\otimes n)\mapsto(c\otimes m)\otimes(1\otimes n)\)。在纯张量上检查两复合为恒等，便得同构。
\end{proof}

有限生成条件在这里不能悄悄删去。例如对 \(R=\mathbb Z_{(p)}\)、\(M=\mathbb Q\)，有 \(\widehat M=0\)，因为 \(p^nM=M\)；但 \(\widehat R\otimes_RM=\mathbb Z_p[1/p]\ne0\)。最后非零可由 \(\mathbb Z_p\) 是整环看出：非零 \(p\) 进数有最低非零数位，相乘时阶数相加，故任何 \(p^n\) 都不是零因子。任意模的完备化因此并不等于沿完备环的普通标量扩张。

\subsection{Noether 局部环完备化的忠实平坦性}
\label{b3:subsec:1304:02}

\begin{theorem}\label{b3:thm:completion-flat}
若 \(R\) 为交换 Noether 环、\(I\subseteq R\) 为理想，则其 \(I\) 进完备化 \(\widehat R\) 是平坦 \(R\) 模。若 \(R\) 还是局部环，极大理想为 \(\mathfrak m\)，且 \(I\subseteq\mathfrak m\)，则 \(R\rightarrow\widehat R\) 忠实平坦。特别地，极大理想进完备化忠实平坦。
\end{theorem}
\begin{proof}
对任意有限生成理想 \(J\subseteq R\)，映射
\[
J\otimes_R\widehat R\longrightarrow R\otimes_R\widehat R=\widehat R
\]
由自然张量表示等同于 \(\widehat J\rightarrow\widehat R\)。\cref{b3:thm:completion-exact}说明后一个映射单射。于是\cref{b3:thm:ideal-flatness-criterion}给出平坦性。

若 \(R\) 局部且 \(I\subseteq\mathfrak m\)，则 \(R/\mathfrak m\) 的 \(I\) 进滤过第一步就为零。由\cref{b3:prop:completion-quotients}，
\[
\widehat R/\mathfrak m\widehat R\cong R/\mathfrak m\ne0.
\]
\(R\) 只有这一个极大理想，故\cref{b3:thm:faithfully-flat-criterion}给出忠实平坦性。
\end{proof}

其直接含义是：一个 \(R\) 模即使不是有限生成的，也不会在张量到 \(\widehat R\) 后从非零变成零；模映射的单射、满射及正合性也能在张量后检测。这里谈的是张量积函子。上一小节的 \(M=\mathbb Q\) 例子恰好提醒我们，对任意模不能把它替换成完备化函子。

\subsection{极大理想、剩余域与维数}
\label{b3:subsec:1304:03}

\begin{proposition}\label{b3:prop:completed-local-ring}
设 \((R,\mathfrak m)\) 为交换 Noether 局部环，\(k=R/\mathfrak m\) 为剩余域，\(B=\widehat R\) 为极大理想进完备化。则 \(B\) 是 Noether 局部环，其极大理想为 \(\mathfrak n=\mathfrak mB\)，剩余域自然等于 \(k\)。对每个 \(r\ge1\) 有
\[
B/\mathfrak n^r\cong R/\mathfrak m^r,
\qquad \operatorname{gr}_{\mathfrak n}(B)\cong\operatorname{gr}_{\mathfrak m}(R).
\]
\end{proposition}
\begin{proof}
Noether 性和两项同构分别由\cref{b3:thm:completion-noetherian}和\cref{b3:prop:completion-quotients}给出。商 \(B/\mathfrak n=k\) 是域，所以 \(\mathfrak n\) 为极大理想。若 \(b\notin\mathfrak n\)，选 \(a\in R\) 使 \(ab\equiv1\pmod{\mathfrak n}\)，并写 \(ab=1-u\)、\(u\in\mathfrak n\)。完备性使 \(\sum_{j\ge0}u^j\) 收敛，逐有限阶等比恒等式给出 \((1-u)\sum_ju^j=1\)。故 \(b\) 可逆，其他元素都不能属于极大理想，\(B\) 为局部环。
\end{proof}

所有有限阶商及剩余域保持不变，并不表示原环与完备环相同。\(k[x]_{(x)}\) 中的元素是原点分母非零的有理函数，而 \(k[[x]]\) 还包含大量并非有理函数的级数。维数保持是另一项需要证明的结论，将由本节末的长度增长比较给出。

\subsection{由完备化正合性到忠实平坦性}
\label{b3:subsec:1304:04}

上面的论证可按每步实际使用的假设回看。Artin–Rees 需要 Noether 环上的有限模，由此把诱导滤过换成子模自身的滤过；逆极限提升随即给出有限模正合性。有限展示和张量右正合性给出 \(\widehat M\cong M\otimes_R\widehat R\)。对有限生成理想取 \(M=J\)，便得到理想张量映射单射，进而推广为对所有模的平坦性。最后剩余域非零才负责“忠实”二字。

\ProseParagraphBreak
全局完备化未必忠实。例如 \(\mathbb Z\rightarrow\mathbb Z_p\) 平坦，但取素数 \(q\ne p\)，\(q\) 在 \(\mathbb Z_p\) 中可逆，所以
\[
(\mathbb Z/q\mathbb Z)\otimes_{\mathbb Z}\mathbb Z_p=0.
\]
这不违反局部结论，因为 \(\mathbb Z\) 不是在 \((p)\) 处的局部环；先局部化为 \(\mathbb Z_{(p)}\) 再完备化，便是忠实平坦映射。忠实性判断的是保留哪些模的信息，不是把“平坦”换成一个更强的形容词而不增加假设。

\subsection{Hilbert–Samuel 理论与完备化的维数保持}
\label{b3:subsec:1304:05}

\begin{theorem}\label{b3:thm:completion-dimension}
对交换 Noether 局部环 \((R,\mathfrak m)\) 及其极大理想进完备化 \(B=\widehat R\)、\(\mathfrak n=\mathfrak mB\)，有
\[
\dim B=\dim R,\qquad e_{\mathfrak n}(B)=e_{\mathfrak m}(R).
\]
更一般地，对非零有限 \(R\) 模 \(M\)，
\[
\dim_B\widehat M=\dim_RM,\qquad
 e_{\mathfrak n}(\widehat M)=e_{\mathfrak m}(M).
\]
\end{theorem}
\begin{proof}
\cref{b3:prop:completion-quotients}给出
\(\widehat M/\mathfrak n^{r+1}\widehat M\cong M/\mathfrak m^{r+1}M\)。两者的标量作用都通过同一个商环 \(B/\mathfrak n^{r+1}\cong R/\mathfrak m^{r+1}\) 分解，所以子模及组成列完全相同，特别是
\[
\ell_B(\widehat M/\mathfrak n^{r+1}\widehat M)
=\ell_R(M/\mathfrak m^{r+1}M).
\]
\(B\) 已证为 Noether 局部环，\(\widehat M\) 已证有限且由忠实平坦性非零，故两侧都能应用\cref{b3:thm:hilbert-samuel}。相等的累计函数给出相等的最终多项式，比较次数与首项即得维数和重数。取 \(M=R\) 得环的结论。
\end{proof}

因此完备化保留局部长度增长及其决定的维数，却未必保留素理想分解的所有细节。例如曲线在形式邻域中可能分出原环中尚未分开的分支；维数相等并不声称两者素谱同胚。
